% Pursuing Stacks (cotes 134-2 to 134-8): archive cote 134-6, site page 39 (PDF page 40), typescript page 357.
% typescript. Draft compared: ps5.tex:3268-3343.
% AI-transcribed by Claude Opus 5.5 (claude-opus-5-5, adaptive thinking, effort high), Message Batches API, batch ps3 (msgbatch_01XdFmuVVBAcgc7EFxdtCqsj), 2026-11-12; not yet checked by eye.
\scanpage{134-6/39}
\mspage{357}

tion when instead of $f\colon A \to B$ , we got a functor
\[ (6) \qquad f : A \longrightarrow B^{\wedge} \quad , \quad \text{factoring through } B^{\wedge}_{c} \xout{\unsure{;}} \]
\xout{orxequival} (which need not factor through B), or equivalently a func-%
tor
\[ (7) \qquad f_{!} : A^{\wedge} \longrightarrow B^{\wedge} \]
commuting with small direct limits, or equivalently still, a functor
$f^{*}$ in opposite direction, commuting with small inverse limits, but
in the last two cases with the extra condition that $f_{!}(A) \subset B^{\wedge}_{c}$ . We
may want to extend the notion of morphism of contractors to include
this situation, hence expressed by the two \inserted{following} conditions on a functor
f (6) \unsure{o\over{f}{r}|of} $f_{!}$ (7) \xout{\unsure{,}}, or on the pair $(f_{!},f^{*})$ of adjoint functors:
\[ (8) \qquad \begin{cases} f_{!}(A) \subset B^{\wedge}_{c} \\ f^{*}(B) \subset A^{\wedge}_{c} \qquad . \end{cases} \]
However, in order for this notion to be stable under composition,
we should strengthen the first of the relations (8) into
\[ (9) \qquad f_{!}(A^{\wedge}_{c}) \subset B^{\wedge}_{c} \quad , \]
which follows automatically whenever $f_{!}$ commutes to finite products
(cf. section 8\over{9}{5}), but may not follow from (8) in general, even in
the case when f factors through B i.e. in the case we start with a
functor (0) $f\colon A \longrightarrow B$. Thus we get \uline{two} plausible notions of \inserted{a} morphism
of contractors, neither of which implies the other, and I feel unable
to predict which one will prove the more useful. As far as terminology
goes, it seems reasonable to reserve the name ``morphism of contractor\hand{s''}
to the first notion, as the second is adequately characterized as a
\uline{bimorphism} between the contractibility structures $(A^{\wedge},A^{\wedge}_{c})$ and $(B^{\wedge},B^{\hand{\wedge}}_{c})$
defined by the contractors A and B\xout{\unsure{,}} (cf. section 86).

16.7.

98. Next point on my provisional program is induced structures
(asphericity of contractibility structures) on a category $M_{/a}$ , when
one is given on M - but finally I decided to skip this, as there was
no urgent need for clarifying this and I am not writing a treatise,
thanks Gods ! I felt more interested writing down the proof that the
category $\hand{\Delta}^{f}$ of standard ordered simplices without degeneracies is a
weak test category, as announced months ago,in section 43. It then

\margin{39}
